\originalchapter{作业 7}
MIT 18.330, Spring 2012. 原件: Homework 7, Due Thursday May 3. 以下解答均为 AI 编写, 非官方答案. 本章使用
$\widehat f(k)=\int_\R f(x)\ee^{-\ii kx}\dd x$,
$f(x)=(2\pi)^{-1}\int_\R\widehat f(k)\ee^{\ii kx}\dd k$ 的约定.

\begin{exercise}\label{ps7-q1}
(原题 1, 3 分) 证明 Fourier 变换的下列性质 ($x,k\in\R$):
\begin{enumerate}[label=(\alph*)]
\item 缩放: 若 $g(x)=f(x/a)$, $a>0$, 则 $\widehat g(k)=a\widehat f(ak)$.
\item 共轭: 若 $g(x)=\overline{f(x)}$, 则 $\widehat g(k)=\overline{\widehat f(-k)}$.
\item 对称性: 若 $f$ 实值且为偶函数 ($f(x)=f(-x)$), 则 $\widehat f$ 也实值且为偶函数.
\item 对称性: 若 $f$ 实值且为奇函数 (原件括号却印为 $f(x)=f(-x)$), 证明 $\widehat f$ 纯虚且为奇函数.
\end{enumerate}
\end{exercise}
\begin{solution}
以 $f\in L^1(\R)$ 为足以保证各积分存在的条件.
(a) 代换 $x=ay$, $\dd x=a\dd y$, 得
$\widehat g(k)=a\int_\R f(y)\ee^{-\ii(ak)y}\dd y=a\widehat f(ak)$.
(b) 对积分取共轭,
$\widehat g(k)=\overline{\int_\R f(x)\ee^{\ii kx}\dd x}=\overline{\widehat f(-k)}$.
(c) 分解为余弦和正弦积分; $f(x)\sin(kx)$ 为奇函数, 其积分为零, 所余余弦积分实值且关于 $k$ 为偶.
(d) 整理注: 按“奇函数”将括号修为 $f(-x)=-f(x)$. 此时 $f(x)\cos(kx)$ 为奇函数, 故余弦积分消失;
$\widehat f(k)=-\ii\int_\R f(x)\sin(kx)\dd x$, 它纯虚且关于 $k$ 为奇. 若同时照抄括号的偶条件与文字的奇条件, 只剩 $f=0$, 不能代表原问意图.
\end{solution}

\begin{exercise}\label{ps7-q2}
(原题 2, 3 分) 定义 $\operatorname{sinc}(x)=\sin x/x$, $x=0$ 处以连续延拓取值 $1$.
\begin{enumerate}[label=(\alph*)]
\item 证明 sinc 不可积, 即不属于 $L^1(\R)$. 提示: 若 $f\ge g\ge0$ 且 $\int g$ 发散, 则 $\int f$ 也发散.
\item 含 sinc 的某些积分仍有意义. 根据 Fourier 变换理论预测 $\int_{-\infty}^{\infty}\operatorname{sinc}(x)\dd x$ 的值.
\end{enumerate}
\end{exercise}
\begin{solution}
(a) 在 $I_n=[n\pi+\pi/6,n\pi+5\pi/6]$ ($n\ge1$) 上,
$|\sin x|\ge1/2$, $x\le(n+1)\pi$, 故
$\int_{I_n}|\sin x/x|\dd x\ge1/[3(n+1)]$.
对互不重叠的这些区间求和得到发散的调和下界, 因而 $\int_\R|\operatorname{sinc}x|\dd x=\infty$.
(b) 区间指示函数 $\mathbf1_{[-1,1]}(k)$ 的逆变换为
$(2\pi)^{-1}\int_{-1}^1\ee^{\ii kx}\dd k=\sin x/(\pi x)$.
因此在广义 Fourier 意义下 $\widehat{\operatorname{sinc}}=\pi\mathbf1_{[-1,1]}$, 零频率对应积分 $\pi$.
普通反常积分也确实存在: 在 $[A,B]$ 对 $\sin x/x$ 分部积分, 尾部绝对值不超过 $3/A$, 从而满足 Cauchy 条件.
为核对常数, 对 $\epsilon>0$ 有
$\int_0^\infty \ee^{-\epsilon x}\sin x/x\dd x=\arctan(1/\epsilon)$: 将 $\sin x$ 换成 $\sin(bx)$ 后对 $b$ 求导, 得 $\epsilon/(\epsilon^2+b^2)$, 从 $b=0$ 积至 $1$.
分部积分给出对 $\epsilon\ge0$ 一致的 $O(1/A)$ 尾界, 故可先控尾再令 $\epsilon\downarrow0$, 得半轴积分为 $\pi/2$. sinc 为偶函数, 双半轴反常积分为 $\pi$, 并非绝对收敛的 Lebesgue 积分.
\end{solution}

\begin{exercise}\label{ps7-q3}
(原题 3, 2 分) 设核 $K(x)$, 考虑积分方程
\[
u(x)+\int_{-\infty}^{\infty}K(x-y)u(y)\dd y=f(x).
\]
类似方程出现于计算机图形渲染以及雷达波被飞机散射等问题. 用 Fourier 变换给出 $u$ 的公式. $K$ 满足何种条件可避免除零?
\end{exercise}
\begin{solution}
积分项为卷积 $K*u$, 变换后
$(1+\widehat K(k))\widehat u(k)=\widehat f(k)$.
只要 $1+\widehat K(k)\ne0$, 形式上
\[
u(x)=\frac1{2\pi}\int_\R
\frac{\widehat f(k)}{1+\widehat K(k)}\ee^{\ii kx}\dd k.
\]
整理注: 点态不为零仅避免逐点除零, 还不足以保证任意 $L^2$ 右端的稳定求解. 一个常用充分条件是 $K\in L^1(\R)$ 且
$\inf_{k\in\R}|1+\widehat K(k)|\ge c>0$; 则该乘子在 $L^2$ 上有界可逆, 公式按 $L^2$ 逆变换理解. 例如 $\|K\|_1<1$ 即足够, 因 $|\widehat K|\le\|K\|_1$. 对更强的点态积分公式, 还需被积函数可积等额外条件.
\end{solution}

\begin{exercise}\label{ps7-q4}
(原题 4, 2 分) 证明 $v_k(x)=c\ee^{-\ii kx}$ ($k\in\Z$) 在 $[0,2\pi]$ 上关于内积
$\langle f,g\rangle=\int_0^{2\pi}f(x)\overline{g(x)}\dd x$ 正交, 并求使其规范化的 $c$.
\end{exercise}
\begin{solution}
对 $k\ne\ell$,
$\langle v_k,v_\ell\rangle=|c|^2[\ee^{\ii(\ell-k)x}/(\ii(\ell-k))]_0^{2\pi}=0$;
对 $k=\ell$ 内积为 $2\pi|c|^2$. 所以通常取正实常数 $c=1/\sqrt{2\pi}$; 任意模为该值的复数也可以. 原件内积中的共轭线与 (b) 的共轭线均按实际 PDF 核对保留, 不按丢失共轭的文本提取结果改题.
\end{solution}
